% preamble.tex -- typesetting for Antifascist Intelligence.
% Hand-edited. This is the design surface: change type, measure, and styling
% here. Nothing regenerates this file.
%
% The document class is deliberately NOT declared here. It is in book.tex,
% above the \input that pulls this file in (D-170). Overleaf picks a project's
% main file by searching the .tex files for that declaration, so while it lived
% here Overleaf compiled preamble.tex, which has no \begin{document}, and
% aborted with "no legal \end found". The class options travel with the master
% for that reason. Do not move them back.
%
% The declaration is spelled out nowhere in this file, comment included, in
% case that search does not skip comments -- which is not checkable from here.

\usepackage{fontspec}
\setmainfont{TeX Gyre Pagella}          % a Palatino, matching the HTML path
\usepackage{microtype}
\usepackage[a4paper,margin=1in]{geometry}
\usepackage{fancyhdr}
\usepackage{tcolorbox}
\usepackage{longtable,array}   % the condensed text's 18 tables, with long cells (D-621)
\tcbuselibrary{breakable}   % the only library the box style needs
\usepackage[english]{babel}
\usepackage{textcomp}
\usepackage{todonotes}
% refs.bib uses \euro, which is not a base LaTeX macro. Everything else it
% uses (\url \S \i \emph \c \v \textsection \L) is standard.
\providecommand{\euro}{\texteuro}
\usepackage[hidelinks]{hyperref}   % before biblatex, per its own guidance
\usepackage[backend=biber,style=authoryear,sorting=nyt,maxcitenames=2]{biblatex}
% Let long URLs in the References break after digits and letters, not only at / and .:
% the OHCHR docstore links are one unbreakable string and ran off the page.
\setcounter{biburlnumpenalty}{100}
\setcounter{biburlucpenalty}{100}
\setcounter{biburllcpenalty}{100}
\addbibresource{../finishing/refs.bib}

% ---------------------------------------------------------------------------
% Draft apparatus: the DRAFT watermark and the status footer (D-319).
%
% One switch governs all of it. Set \draftmodefalse to ship a clean book;
% nothing else in this file has to change, and the two generated counts below
% go unread rather than stale.
\newif\ifdraftmode \draftmodetrue

% The two numbers the footer prints, and the list of entries a human has checked
% (one \refcheckedkey{KEY} per entry), which the citation highlight below reads.
% GENERATED by finishing/tools/refs_ledger.py from finishing/refs-ledger.tsv and
% verified current by check_all.sh, so a page cannot print a count or a check
% the ledger does not carry.
\newcommand{\refcheckedkey}[1]{\csdef{refok@#1}{}}
%% draft-status.tex -- GENERATED by finishing/tools/refs_ledger.py from
%% finishing/refs-ledger.tsv. Do not edit by hand; edits here are overwritten.
%%
%% The footer on every page prints these two numbers, and the draft highlight
%% leaves unmarked every citation of a key listed below by \refcheckedkey. They
%% are regenerated from the ledger and checked by check_all.sh, so the page cannot
%% claim a count or a check the ledger does not carry.
\newcommand{\refschecked}{2}
\newcommand{\refstotal}{575}
\refcheckedkey{abraham2024lavender}
\refcheckedkey{rickandmorty2023}


% The build timestamp, MM/DD/YYYY HH:MM, computed by TeX as the job starts.
% Deliberately NOT written into draft-status.tex: a generated file that changed
% on every build would make that file's --check meaningless, and the check is
% the whole reason the footer can be trusted. Pure TeX arithmetic rather than
% \directlua so that the Overleaf round trip (D-169) still compiles if Overleaf
% picks an engine that is not LuaTeX.
\makeatletter
\newcount\ds@hh \newcount\ds@mm \newcount\ds@tmp
\newcommand{\ds@pad}[1]{\ifnum#1<10 0\fi\number#1}
\AtBeginDocument{%
  \ds@hh=\time \divide\ds@hh by 60
  \ds@mm=\time \ds@tmp=\ds@hh \multiply\ds@tmp by 60 \advance\ds@mm by -\ds@tmp
  \edef\buildstamp{%
    \ds@pad{\month}/\ds@pad{\day}/\number\year\space\ds@pad{\ds@hh}:\ds@pad{\ds@mm}}%
}
\makeatother

% The separator is a middot rather than the hyphen the line was sketched with:
% the book's own dashes are all spaced em dashes and a run of those in a
% four-part footer reads as heavily as the prose it sits under.
\newcommand{\draftstatusline}{%
  PREPRINT \textperiodcentered\ WORKING MANUSCRIPT \textperiodcentered\
  Rendered \buildstamp\ \textperiodcentered\
  \refschecked/\refstotal\ bibliographic entries human-checked}

\ifdraftmode
  \usepackage{xcolor}          % loaded by tcolorbox above; named here because
                               % \draftstatusline's grey depends on it
  \usepackage{draftwatermark}
  % Angle, size and position are the package's own defaults -- 45 degrees,
  % 0.25\paperwidth, centred -- and are left unset rather than restated. Only
  % the lightness moves: the default 0.8 grey is dark enough to read through on
  % a page of 11pt Pagella, and 0.88 stays legible as a watermark without
  % competing with the prose it sits behind.
  \SetWatermarkText{DRAFT}
  \SetWatermarkLightness{0.88}

  % Unchecked citations (D-641). Every in-text citation of an entry the ledger
  % does not mark checked is set on an orange highlight; a checked entry's
  % citations are set plain. The test is per key, inside biblatex's per-key
  % `cite' macro, so in \autocites{a}{b} each key carries its own state and the
  % brackets and pre/postnotes around them stay unmarked. To clear one, a human
  % marks it: finishing/tools/refs_ledger.py --mark KEY --by INITIALS (or
  % --toggle), then rebuild.
  %
  % Three renderings of the one mark. Under tex4ht (the HTML proof) a span the
  % page's stylesheet colours, since tex4ht can't carry lua-ul's highlight.
  % Under LuaTeX, lua-ul's \highLight, which breaks across lines and keeps the
  % citation's link working. Under any other engine -- an Overleaf compile on
  % XeLaTeX, say (D-169) -- orange text, because lua-ul needs LuaTeX.
  \ifdefined\HCode
    \newcommand{\refunchecked}[1]{%
      \HCode{<span class="ref-unchecked">}#1\HCode{</span>}}
  \else\ifdefined\directlua
    \usepackage{luacolor,lua-ul}
    \newcommand{\refunchecked}[1]{\highLight[orange!45]{#1}}
  \else
    \newcommand{\refunchecked}[1]{\textcolor{orange!80!black}{#1}}
  \fi\fi
  \letbibmacro{draft:cite}{cite}
  \renewbibmacro*{cite}{%
    \ifcsdef{refok@\thefield{entrykey}}
      {\usebibmacro{draft:cite}}
      {\refunchecked{\usebibmacro{draft:cite}}}}
\fi

\setcounter{secnumdepth}{2}             % deepest heading in the book is x.y.z
% The contents list chapters and sections and stops there. At depth 2 it ran
% pages 2-5 of a 196-page book, three full pages and an 11-line stub, and the 64
% subsections were most of it. What the depth was buying is a lookup path for the
% 79 cross-references that point at an x.y.z -- 35 percent of the 228 -- and the
% running head below now carries that instead, on the page the reader is already
% on. Cutting the depth without the head first would have left those 79 with no
% lookup at all, the class being oneside (D-260).
\setcounter{tocdepth}{1}

% Run-in head: a bold lead-in line above a paragraph, used 169 times.
\newcommand{\runin}[1]{\par\medskip\noindent\textbf{#1}\par\nopagebreak\smallskip}

% A one-line standing note under a heading, saying which of the preface's
% five claims the text beneath it rests on. Small italic, so it reads as a
% note and not as the first paragraph.
\newcommand{\standing}[1]{\noindent{\small\itshape #1}\par\medskip}

\newtcolorbox{esbox}{
  colback=black!3, colframe=black!35, boxrule=0.4pt, arc=1pt,
  left=8pt, right=8pt, top=8pt, bottom=8pt, breakable,
  fontupper=\small,
  before upper={\setlength{\parskip}{0.5\baselineskip}%
               \setlength{\parindent}{0pt}},
}
\newcommand{\boxtitle}[1]{{\bfseries #1}\par\smallskip}

% Front and back matter are \chapter*, which does not set \@currentlabel, so a
% plain \label after one captures whatever number was last set -- silently
% wrong. This pins the value a \ref to the heading yields, which is the
% ordinal of the heading and not a number the page shows.
\makeatletter
\newcommand{\unnumberedlabel}[2]{\begingroup\def\@currentlabel{#2}\label{#1}\endgroup}
\makeatother

% \chapter* also sets no running mark, so the head left standing by the last
% numbered chapter persists over the back matter -- every page of the glossary
% read "CHAPTER 12. CONCLUSION AND OUTLOOK" (Q-055). Called once per starred
% chapter, after the heading.
\newcommand{\backmattermark}[1]{\markboth{\MakeUppercase{#1}}{\MakeUppercase{#1}}}

% The running head carries the deepest heading in force, so a reader who follows
% a cross-reference to an x.y.z can find it from the page they are on. The class
% is oneside, where book.cls drops \sectionmark entirely and puts the chapter on
% every page -- all 196 of them read "CHAPTER n. TITLE" and told a reader nothing
% they did not know. \subsectionmark is defined because \@sect calls it and the
% standard page styles \@gobble it, which is why the subsection never reached a
% head. Small caps rather than the class's uppercase: an x.y.z title set in full
% capitals is long and hard to read, and \backmattermark's own uppercase marks
% still set as capitals here, so the two agree (D-260).
\pagestyle{fancy}
\fancyhf{}
\fancyhead[L]{\small\scshape\rightmark}
\fancyhead[R]{\small\thepage}
\ifdraftmode\fancyfoot[C]{\scriptsize\color{black!55}\draftstatusline}\fi
\renewcommand{\headrulewidth}{0pt}
\renewcommand{\chaptermark}[1]{\markboth{\MakeUppercase{#1}}{\thechapter.\ #1}}
\renewcommand{\sectionmark}[1]{\markright{\thesection.\ #1}}
\renewcommand{\subsectionmark}[1]{\markright{\thesubsection.\ #1}}
% A chapter's opening page is \thispagestyle{plain}; it carries the folio alone.
\fancypagestyle{plain}{%
  \fancyhf{}%
  \fancyhead[R]{\small\thepage}%
  \ifdraftmode\fancyfoot[C]{\scriptsize\color{black!55}\draftstatusline}\fi%
  \renewcommand{\headrulewidth}{0pt}%
}

% The title page. book.tex supplies \booktitlemain, \booksubtitle and
% \bookauthor, and calls \booktitlepage as the first thing after \frontmatter.
% The book had no title page at all until 2026-08-25 (D-073); it opened on the
% table of contents, so its own title appeared nowhere in the PDF.
\newcommand{\booktitlepage}{%
  \begin{titlepage}
    \centering
    \vspace*{0.28\textheight}
    {\fontsize{30}{36}\selectfont \booktitlemain\par}
    \vspace{1.2\baselineskip}
    {\Large\itshape \booksubtitle\par}
    \vspace{3\baselineskip}
    {\large \bookauthor\par}
    \vfill
    {\footnotesize
     Copyright \textcopyright{} 2026 \bookauthor.\\
     Licensed under CC BY-NC-ND 4.0.\\[0.4\baselineskip]
     Third-party material quoted in this work, including the epigraphs,\\
     remains the property of its rights holders and is reproduced under\\
     fair use. The license above does not extend to it.\par}
    \ifdraftmode
      \vspace{1.2\baselineskip}
      {\footnotesize\color{black!55}\draftstatusline\par}
    \fi
    \vspace{0.06\textheight}
  \end{titlepage}%
}
